\documentclass[12pt,a4paper]{article}
\usepackage[utf8]{inputenc}
\usepackage[T2A]{fontenc}
\usepackage[russian]{babel}
\usepackage{amsmath,amssymb}
\usepackage[left=25mm,right=25mm,top=20mm,bottom=20mm]{geometry}
\usepackage{enumitem}

\newcommand{\dd}{\mathrm{d}}
\newcommand{\pp}{\partial}
\newcommand{\ext}{\mathrm{ext}}
\newcommand{\eff}{\mathrm{eff}}
\newcommand{\st}{\text{ст}}
\newcommand{\sat}{\mathrm{sat}}

\title{Метод интегральных соотношений (моментов)\\
для задачи двухфазной фильтрации Black~Oil\\
с балансом энергии в радиальной постановке}
\author{}
\date{}

\begin{document}
\maketitle
\tableofcontents
\newpage

%=============================================================================
\section{Введение}\label{sec:intro}
%=============================================================================

Полная система уравнений двухфазной фильтрации Black~Oil с балансом
энергии состоит из трёх ДУвЧП (дифференциальных уравнений в частных производных)
относительно функций $p(r,t)$, $T(r,t)$, $S_g(r,t)$ в области $r\in[r_w, R_\ext]$,
$t \geqslant 0$.

Предлагается следующая стратегия приближённого решения задачи. Метод интегральных соотношений
(моментов) Баренблатта~\cite{barenblatt1954} применяется ко
\emph{всем трём уравнениям}. ДУвЧП сводятся к системе ОДУ (обыкновенных дифференциальных уравнений) 
для параметров профилей $p(r,t)$, $S_g(r,t)$ и $T(r,t)$.
Температурное поле входит в уравнения давления и
насыщенности через PVT-свойства $R_s(p,T)$, $B_o(p,T)$,
$\mu_o(p,T)$ и т.\,д., а также через члены уравнений
$C_T\,\pp T/\pp t$ и $\beta_o^\sat\,\pp T/\pp t$. Таким образом, гидродинамика и температурная задача
оказываются сильно связанными, и <<последовательное>> решение
(сначала гидродинамика, потом температура) может приводить
к неконтролируемым ошибкам.


%=============================================================================
\section{Исходные уравнения}\label{sec:equations}
%=============================================================================

Система состоит из трёх уравнений на три неизвестных
$p(r,t)$, $T(r,t)$, $S_g(r,t)$. Все PVT-свойства зависят
от $p$ и $T$: $R_s(p,T)$, $B_{\alpha}(p,T)$,
$\mu_\alpha(p,T)$, $\alpha = o, g$ и т.\,д.

\subsection{Уравнение давления}

Уравнение для давления принимает следующий вид:
\begin{equation}\label{eq:p_eq}
  \boxed{
    C_p\,\frac{\pp p}{\pp t}
    + C_T\,\frac{\pp T}{\pp t}
    + \mathcal{D} = 0,
  }
\end{equation}
где:

\paragraph{Подвижности фаз:}
\begin{equation}\label{eq:mobility}
  \lambda_\alpha = \frac{k\,k_{r\alpha}(S_g)}{\mu_\alpha(p,T)},
  \qquad \alpha \in \{o, g\}.
\end{equation}

\paragraph{Суммарная приведённая подвижность:}
\begin{equation}\label{eq:Lambda}
  \Lambda(p,T,S_g) = \frac{\lambda_o + \lambda_g}{B_g}
  = \frac{k}{B_g(p,T)}\!\left(
      \frac{k_{ro}(S_g)}{\mu_o(p,T)}
      + \frac{k_{rg}(S_g)}{\mu_g(p,T)}
    \right).
\end{equation}

\paragraph{Коэффициент ёмкости при $\pp p/\pp t$:}
\begin{equation}\label{eq:Cp}
  C_p = \phi\!\left[
    S_o\,(c_r + c_o^\sat)\!\left(\frac{1}{B_g} - \frac{R_s}{B_o}\right)
    + \xi_o\!\left(\frac{\pp R_s}{\pp p}
      + R_s\,(c_r + c_o^\sat)\right)
    + \frac{S_g}{B_g}\,(c_g + c_r)
  \right],
\end{equation}
где $\xi_o = S_o/B_o$, $S_o = 1 - S_g$, а <<полная>> сжимаемость
нефтяной фазы вдоль кривой насыщения:
\begin{equation}\label{eq:co_sat}
  c_o^\sat = c_o
    - \frac{1}{B_o}\frac{\pp B_o}{\pp R_s}\bigg|_{p,T}
      \frac{\pp R_s}{\pp p}\bigg|_T,
  \qquad
  c_o = -\frac{1}{B_o}\frac{\pp B_o}{\pp p}\bigg|_{T,\,R_s}.
\end{equation}

\paragraph{Коэффициент при $\pp T/\pp t$:}
\begin{equation}\label{eq:CT}
  C_T = \phi\!\left[
    -S_o\,\beta_o^\sat\!\left(\frac{1}{B_g} - \frac{R_s}{B_o}\right)
    + \xi_o\!\left(\frac{\pp R_s}{\pp T}
      - R_s\,\beta_o^\sat\right)
    - \frac{S_g}{B_g}\,\beta_g
  \right],
\end{equation}
где
\begin{equation}\label{eq:beta_o_sat}
  \beta_o^\sat = \beta_o
    + \frac{1}{B_o}\frac{\pp B_o}{\pp R_s}\bigg|_{p,T}
      \frac{\pp R_s}{\pp T}\bigg|_p,
  \qquad
  \beta_o = \frac{1}{B_o}\frac{\pp B_o}{\pp T}\bigg|_{p,\,R_s}.
\end{equation}

\paragraph{Пространственный оператор:}
\begin{equation}\label{eq:D_operator}
  \mathcal{D} = -\frac{1}{r}\,\frac{\pp}{\pp r}\!\left(
      r\,\Lambda\,\frac{\pp p}{\pp r}
    \right)
    - \frac{\lambda_o}{B_o}\!\left(
      \frac{\pp R_s}{\pp p}\,\frac{\pp p}{\pp r}
      + \frac{\pp R_s}{\pp T}\,\frac{\pp T}{\pp r}
    \right)\frac{\pp p}{\pp r}.
\end{equation}

\subsection{Уравнение насыщенности}

Уравнение для определения насыщенности принимает следующий вид:
\begin{equation}\label{eq:Sg_eq}
  \boxed{
    \frac{\pp S_g}{\pp t}
    = S_o\,(c_r + c_o^\sat)\,\frac{\pp p}{\pp t}
      - S_o\,\beta_o^\sat\,\frac{\pp T}{\pp t}
      + \frac{B_o}{\phi}\,\frac{1}{r}\,\frac{\pp}{\pp r}\!
        \left(r\,\frac{u_o}{B_o}\right).
  }
\end{equation}
Подставляя $u_o = -\lambda_o\,\pp p/\pp r$, получим
\begin{equation}\label{eq:Sg_eq2}
  \frac{\pp S_g}{\pp t}
  = S_o\,(c_r + c_o^\sat)\,\frac{\pp p}{\pp t}
    - S_o\,\beta_o^\sat\,\frac{\pp T}{\pp t}
    - \frac{B_o}{\phi}\,\frac{1}{r}\,\frac{\pp}{\pp r}\!
      \left(\frac{r\,\lambda_o}{B_o}\,\frac{\pp p}{\pp r}\right).
\end{equation}

\subsection{Уравнение энергии}

Также запишем уравнение баланса энергии в терминах температуры:
\begin{equation}\label{eq:energy}
  \boxed{
    (\rho c_p)_\eff\,\frac{\pp T}{\pp t}
    + (\rho c_p u)_\mathrm{conv}\,\frac{\pp T}{\pp r}
    = \frac{1}{r}\,\frac{\pp}{\pp r}\!\left(
        r\,\lambda_\eff\,\frac{\pp T}{\pp r}
      \right)
    + Q_\mathrm{JT} + Q_\mathrm{ad} + Q_\mathrm{deg},
  }
\end{equation}
где:

\paragraph{Эффективная теплоёмкость:}
\begin{equation}\label{eq:rhocp_eff}
  (\rho c_p)_\eff
  = (1-\phi)\,(\rho c_p)_s
    + \phi\,(S_o\,\rho_o\,c_{p,o} + S_g\,\rho_g\,c_{p,g}).
\end{equation}

\paragraph{Конвективный коэффициент:}
\begin{equation}\label{eq:rhocp_conv}
  (\rho c_p u)_\mathrm{conv}
  = \rho_o\,c_{p,o}\,u_o + \rho_g\,c_{p,g}\,u_g.
\end{equation}

\paragraph{Эффект Джоуля--Томсона:}
\begin{equation}\label{eq:QJT}
  Q_\mathrm{JT} = -\sum_\alpha u_\alpha\,(1 - \alpha_\alpha T)\,
    \frac{\pp p}{\pp r},
\end{equation}
где $\alpha_\alpha$ --- коэффициент теплового расширения
фазы~$\alpha$ при фиксированном составе.

\paragraph{Адиабатический эффект:}
\begin{equation}\label{eq:Qad}
  Q_\mathrm{ad} = \phi\!\left[
    T\sum_\alpha S_\alpha\,\alpha_\alpha
    + c_r\!\left(p + (\rho c_p)_s\,T\right)
  \right]\frac{\pp p}{\pp t}.
\end{equation}

\paragraph{Теплота дегазации:}
\begin{equation}\label{eq:Qdeg}
  Q_\mathrm{deg} = -L(p,T)\,\dot{m},
\end{equation}
где $L(p,T)$ --- удельная теплота дегазации, а массовый темп
выделения газа:
\begin{equation}\label{eq:mdot}
  \dot{m} = -\rho_g^\st\!\left[
    \phi\,\xi_o\!\left(
      \frac{\pp R_s}{\pp p}\,\frac{\pp p}{\pp t}
      + \frac{\pp R_s}{\pp T}\,\frac{\pp T}{\pp t}
    \right)
    + \frac{u_o}{B_o}\!\left(
      \frac{\pp R_s}{\pp p}\,\frac{\pp p}{\pp r}
      + \frac{\pp R_s}{\pp T}\,\frac{\pp T}{\pp r}
    \right)
  \right].
\end{equation}


\subsection{Граничные и начальные условия}

\begin{equation}\label{eq:bc_p}
  p(r_w, t) = p_{wf}(t), \qquad
  p(R_\ext, t) = p_0.
\end{equation}
\begin{equation}\label{eq:bc_T}
  T(R_\ext, t) = T_0, \qquad
  \frac{\pp T}{\pp r}\bigg|_{r = r_w} = 0.
\end{equation}
\begin{equation}\label{eq:ic}
  p(r, 0) = p_0, \qquad T(r, 0) = T_0, \qquad S_g(r, 0) = 0.
\end{equation}

При $t = 0$ давление и температура однородны, газ полностью
растворён ($p_0 > p_b(T_0)$). При снижении $p_{wf} < p_b$
вблизи скважины возникает зона разгазирования
$r \in [r_w, r_b(t)]$, где $p < p_b(T)$.


%=============================================================================
\section{Определение интегральных моментов}\label{sec:moments_def}
%=============================================================================

Перейдем к описанию метода моментов применительно к сформулированной выше 
задаче~\eqref{eq:p_eq}, \eqref{eq:Sg_eq}, \eqref{eq:energy}, \eqref{eq:bc_p}--\eqref{eq:ic}.

\subsection{Зона разгазирования}
Явно выделим подвижную границу (фронт), разделяющую внутреннюю область двухфазного течения 
и внешнюю однофазного течения. Назовем положение этой границы радиусом зоны 
разгазирования $r_b(t)$. Он определяется следующим образом
\begin{equation}\label{eq:rb}
  p(r_b(t), t) = p_b(T(r_b(t), t)).
\end{equation}
При $r > r_b(t)$: $p > p_b(T)$, $S_g = 0$, фильтрация однофазная.

Далее определим зависящие от времени моменты искомых величин.
Моменты "агрегируют" информацию о пространственной зависимости искомых величин,
при этом остаются функциями времени. С учетом радиальной симметрии 
пробные функции выбираются специальным образом, как показано далее.


\subsection{Моменты давления}

\begin{equation}\label{eq:Mn}
  M_n(t) = \int_{r_w}^{R_\ext} r^{2n+1}\,p(r,t)\,\dd r,
  \qquad n = 0, 1, 2, \ldots
\end{equation}

\paragraph{Нулевой момент ($n = 0$):}
\begin{equation}\label{eq:M0}
  M_0(t) = \int_{r_w}^{R_\ext} r\,p(r,t)\,\dd r.
\end{equation}

\paragraph{Первый момент ($n = 1$):}
\begin{equation}\label{eq:M1}
  M_1(t) = \int_{r_w}^{R_\ext} r^3\,p(r,t)\,\dd r.
\end{equation}

\subsection{Момент насыщенности}

Газонасыщенность (как объем отдельной газовой фазы) 
имеет ограниченный носитель $[r_w, r_b(t)]$:
\begin{equation}\label{eq:Sn}
  \mathcal{S}_n(t) = \int_{r_w}^{r_b(t)} r^{2n+1}\,S_g(r,t)\,\dd r.
\end{equation}

$\mathcal{S}_0$ пропорционален объёму свободного газа
в зоне разгазирования.

\subsection{Момент температуры}

Температурное поле определено на всей области $[r_w, R_\ext]$.
Нулевой момент:
\begin{equation}\label{eq:Tn}
  \mathcal{T}_0(t) = \int_{r_w}^{R_\ext} r\,T(r,t)\,\dd r.
\end{equation}

Физический смысл: $2\pi h\,\mathcal{T}_0$ пропорционален
полной тепловой энергии (отклонению от $T_0$) в пласте.


%=============================================================================
\section{Интегральные соотношения}\label{sec:integral_rel}
%=============================================================================

\subsection{Нулевой момент уравнения давления}

Умножим~\eqref{eq:p_eq} на $r$ и проинтегрируем по $r$
от $r_w$ до $R_\ext$. Используя $\mathcal{D}$
из~\eqref{eq:D_operator}:
\begin{multline}\label{eq:moment0}
    \int_{r_w}^{R_\ext} r\,C_p\,\frac{\pp p}{\pp t}\,\dd r
    + \int_{r_w}^{R_\ext} r\,C_T\,\frac{\pp T}{\pp t}\,\dd r
   \\
    = \left[r\,\Lambda\,\frac{\pp p}{\pp r}\right]_{r_w}^{R_\ext}
    + \int_{r_w}^{R_\ext}
      r\,\frac{\lambda_o}{B_o}\!\left(
        \frac{\pp R_s}{\pp p}\,\frac{\pp p}{\pp r}
        + \frac{\pp R_s}{\pp T}\,\frac{\pp T}{\pp r}
      \right)\frac{\pp p}{\pp r}\,\dd r.
\end{multline}



\subsection{Первый момент уравнения давления}

Умножим~\eqref{eq:p_eq} на $r^3$ и проинтегрируем:
\begin{multline}\label{eq:moment1}
  \int_{r_w}^{R_\ext} r^3\,C_p\,\frac{\pp p}{\pp t}\,\dd r
  + \int_{r_w}^{R_\ext} r^3\,C_T\,\frac{\pp T}{\pp t}\,\dd r \\
  = \left[r^3\,\Lambda\,\frac{\pp p}{\pp r}\right]_{r_w}^{R_\ext}
    - 2\int_{r_w}^{R_\ext} r^2\,\Lambda\,\frac{\pp p}{\pp r}\,\dd r
  + \int_{r_w}^{R_\ext}
    r^3\,\frac{\lambda_o}{B_o}\!\left(
      \frac{\pp R_s}{\pp p}\,\frac{\pp p}{\pp r}
      + \frac{\pp R_s}{\pp T}\,\frac{\pp T}{\pp r}
    \right)\frac{\pp p}{\pp r}\,\dd r.
\end{multline}

\subsection{Нулевой момент уравнения насыщенности}

Умножим~\eqref{eq:Sg_eq2} на $r$ и проинтегрируем от $r_w$ до $r_b$:
\begin{multline}\label{eq:Sg_moment0}
  \frac{\dd\mathcal{S}_0}{\dd t}
  = \int_{r_w}^{r_b} r\,S_o\,(c_r+c_o^\sat)\,
      \frac{\pp p}{\pp t}\,\dd r
    - \int_{r_w}^{r_b} r\,S_o\,\beta_o^\sat\,
      \frac{\pp T}{\pp t}\,\dd r 
  - \frac{B_o}{\phi}\!\left[
      \frac{r\,\lambda_o}{B_o}\,\frac{\pp p}{\pp r}
    \right]_{r_w}^{r_b}.
\end{multline}

Член с $\dot{r}_b$ обращается в нуль ($S_g(r_b,t) = 0$,
правило Лейбница). Граничные потоки:
\begin{equation}\label{eq:Sg_flux}
  \left[\frac{r\,\lambda_o}{B_o}\,\frac{\pp p}{\pp r}\right]_{r_w}^{r_b}
  = \frac{r_b\,\lambda_o(r_b)}{B_o(p_b,T_b)}\,p'_b
    - \frac{r_w\,\lambda_o(r_w)}{B_o(p_{wf},T_w)}\,p'_w,
\end{equation}
где $T_b = T(r_b,t)$, $T_w = T(r_w,t)$.


\subsection{Нулевой момент уравнения энергии}

Умножим~\eqref{eq:energy} на $r$ и проинтегрируем
от $r_w$ до $R_\ext$:
\begin{multline}\label{eq:T_moment0}
  \int_{r_w}^{R_\ext} r\,(\rho c_p)_\eff^*\,
    \frac{\pp T}{\pp t}\,\dd r
  + \int_{r_w}^{R_\ext} r\,(\rho c_p u)_\mathrm{conv}^*\,
    \frac{\pp T}{\pp r}\,\dd r \\
  = \left[r\,\lambda_\eff\,\frac{\pp T}{\pp r}
    \right]_{r_w}^{R_\ext}
  + \int_{r_w}^{R_\ext} r\,(Q_\mathrm{JT}^* + Q_\mathrm{ad}^*)\,\dd r.
\end{multline}

Дивергентная часть вычисляется точно:
\begin{equation}\label{eq:T_flux}
  \left[r\,\lambda_\eff\,\frac{\pp T}{\pp r}
    \right]_{r_w}^{R_\ext}
  = R_\ext\,\lambda_\eff\,T'_\ext
    - \underbrace{r_w\,\lambda_\eff\,T'_w}_{= 0},
\end{equation}
поскольку граничное условие~\eqref{eq:bc_T} даёт $T'_w = 0$.


%=============================================================================
\section{Параметризация профилей}\label{sec:profiles}
%=============================================================================

Далее необходимо вычислить интегралы, стоящие в уравнениях, 
определяющих моменты уравнений. Вычисление интегралов опирается на 
задание профилей исхомых функций, их явное параметризованное представление
через функции пространства и времени. Для задания пространственной части учтем,
что газонасыщенность определяется по-разному в двух зонах.


\subsection{Структура области}

Вся расчетная область $[r_w, R_\ext]$ делится на две зоны:
\begin{itemize}[nosep]
  \item \textbf{Зона разгазирования} $r \in [r_w, r_b(t)]$:
        $p < p_b(T)$, $S_g > 0$, двухфазная фильтрация.
  \item \textbf{Однофазная зона} $r \in [r_b(t), R_\ext]$:
        $p \geqslant p_b(T)$, $S_g = 0$, однофазная фильтрация нефти.
\end{itemize}

\subsection{Профиль давления}

\paragraph{Однофазная зона ($r_b \leqslant r \leqslant R_\ext$).} 
В однофазной зоне, вдали от скважины, примем
квазистационарный логарифмический профиль:
\begin{equation}\label{eq:p_outer}
  p(r,t) = p_b^* + (p_0 - p_b^*)\,
    \frac{\ln(r/r_b)}{\ln(R_\ext/r_b)},
  \qquad r_b \leqslant r \leqslant R_\ext.
\end{equation}
Здесь $p_b^* = p_b(T(r_b))$ -- давление насыщения
при температуре на фронте разгазирования.

Альтернатива: профиль можно записать через псевдофункции давления для однофазной фильтрации.
Зесь может получиться хорошая "свадьба" двух подходов.

\paragraph{Зона разгазирования ($r_w \leqslant r \leqslant r_b$).}
В зоне разгазирвоания профиль существенно нестационарный. 
Это должно явно учитываться при параметрзации.
Можно выбрать степенной профиль:
\begin{equation}\label{eq:p_inner}
  p(r,t) = p_{wf} + (p_b^* - p_{wf})\,\Phi\!\left(
    \frac{\ln(r/r_w)}{\ln(r_b/r_w)}
  \right),
\end{equation}
\begin{equation}\label{eq:Phi}
  \Phi(\eta) = \eta^\gamma, \qquad
  \eta = \frac{\ln(r/r_w)}{\ln(r_b/r_w)},
  \qquad \gamma > 0.
\end{equation}
Здесь $p|_{r=r_w}=p_{wf}$ -- давление на стенке скважины.

\paragraph{Непрерывность потока на $r = r_b$.}
При $S_g(r_b) = 0$ подвижности совпадают. 
Тогда производные уравнений~\eqref{eq:p_outer} и~\eqref{eq:p_inner} совпадают. 
Определим перепады давлений в двух зонах $\Delta p_1 = p_b^* - p_{wf}$,
$\Delta p_2 = p_0 - p_b^*$, и из равенства скоростей фильтрации 
в точке $r=r_b(t)$ получим выражение для степени $\gamma$
\begin{equation}\label{eq:gamma_constraint}
  \gamma(t) = \frac{\Delta p_2(t)\,\ln(r_b(t)/r_w)}{\Delta p_1(t)\,\ln(R_\ext/r_b(t))}.
\end{equation}
Величина $p_b^* = p_b(T(r_b))$ зависит от времени через функцию $\Delta T_w(t)$
в профиле температуры~\eqref{eq:T_profile}:
\begin{equation}\label{eq:pb_star}
  p_b^* = p_b\!\left(T_0 - \Delta T_w\,\Psi(\zeta_b)\right),
  \qquad
  \zeta_b = \frac{\ln(R_\ext/r_b(t))}{\ln(R_\ext/r_w)} = \frac{L_e(t)}{L_0}.
\end{equation}

% Производные давления:
% \begin{equation}\label{eq:dp_inner}
%   \frac{\pp p}{\pp r}\bigg|_\text{inner}
%   = \frac{\Delta p_1\,\gamma}{r\,L_b}\,\eta^{\gamma-1},
% \end{equation}
% \begin{equation}\label{eq:dp_outer}
%   \frac{\pp p}{\pp r}\bigg|_\text{outer}
%   = \frac{\Delta p_2}{r\,L_e}.
% \end{equation}

% \paragraph{Обозначения.}
% $L_b = \ln(r_b/r_w)$, $L_e = \ln(R_\ext/r_b)$,
% $\Delta p_1 = p_b^* - p_{wf}$, $\Delta p_2 = p_0 - p_b^*$
% (здесь и далее $p_b^* = p_b(T(r_b))$, см.~\eqref{eq:pb_star}).


\subsection{Профиль насыщенности}

Газонасыщенность имеет компактный носитель $[r_w, r_b(t)]$:
\begin{equation}\label{eq:Sg_profile}
  S_g(r,t) = S_{g,w}(t)\!\left[
    1 - \left(\frac{\ln(r/r_w)}{\ln(r_b/r_w)}\right)^{\!\beta}
  \right],
  \qquad r_w \leqslant r \leqslant r_b,
\end{equation}
$S_g(r_w,t) = S_{g,w}(t)$, $S_g(r_b) = 0$, $\beta > 0$.

Этот профиль также можно параметризовать по-разному.


\subsection{Профиль температуры}\label{ssec:T_profile}

Температурное поле определено на $[r_w, R_\ext]$ и удовлетворяет граничным условиям:
\begin{itemize}[nosep]
  \item $T(R_\ext, t) = T_0$ (контур питания);
  \item $T'(r_w, t) = 0$ (адиабатическая стенка).
\end{itemize}

Профиль параметризуем функцией формы
с амплитудным параметром $\Delta T_w(t) = T_0 - T_w(t)$,
где $T_w(t) = T(r_w, t)$ --- температура на забое:
\begin{equation}\label{eq:T_profile}
  T(r,t) = T_0 - \Delta T_w(t)\,\Psi(\zeta),
  \qquad
  \zeta = \frac{\ln(R_\ext/r)}{\ln(R_\ext/r_w)} = \frac{\ln(R_\ext/r)}{L_0}.
\end{equation}

Функция формы $\Psi(\zeta)$ должна удовлетворять:
$\Psi(0) = 0$ (контур питания),
$\Psi(1) = 1$ (забой),
$\Psi'(1) = 0$ (адиабатическая стенка: $\pp T/\pp r|_{r_w} = 0$).

Простейший профиль с этими свойствами:
\begin{equation}\label{eq:Psi}
  \boxed{
    \Psi(\zeta) = \zeta^\sigma(2 - \zeta^\sigma),
    \qquad \sigma > 0.
  }
\end{equation}

Проверка:
\begin{itemize}[nosep]
  \item $\Psi(0) = 0$. \checkmark
  \item $\Psi(1) = 1\cdot(2-1) = 1$. \checkmark
  \item $\Psi'(\zeta) = 2\sigma\zeta^{\sigma-1}(1 - \zeta^\sigma)$,
        поэтому $\Psi'(1) = 2\sigma\cdot 1\cdot 0 = 0$. \checkmark
\end{itemize}

Производная температуры:
\begin{equation}\label{eq:dT_dr}
  \frac{\pp T}{\pp r}
  = \frac{\Delta T_w}{r\,L_0}\,\Psi'(\zeta)
  = \frac{2\sigma\,\Delta T_w}{r\,L_0}\,
    \zeta^{\sigma-1}(1-\zeta^\sigma).
\end{equation}
При $r = r_w$ ($\zeta = 1$): $\pp T/\pp r = 0$. \checkmark
Максимум $|\pp T/\pp r|$ достигается внутри области.

\paragraph{Интегралы.}
Моменты профиля выражаются через:
\begin{equation}\label{eq:IPsi}
  J^{(m)} = \int_0^1 \zeta^m\,(2 - \zeta^\sigma)\,\dd\zeta
  = \frac{2}{m+1} - \frac{1}{m+\sigma+1}.
\end{equation}

\paragraph{Выбор $\sigma$.}
При преобладании конвекции ($\mathrm{Pe} \gg 1$)
нефтяной поток подводит невозмущённую
температуру $T_0$ от контура, а охлаждение
генерируется локально (JT-эффект, дегазация).
Возмущение $\Delta T$ концентрируется
вблизи $r_w$. При $\sigma = 1$: $\Psi = \zeta(2-\zeta) = 1-(1-\zeta)^2$ ---
параболический профиль, основное изменение
вблизи забоя. При $\sigma > 1$ профиль
ещё сильнее локализован у $r_w$.

$\sigma$ можно зафиксировать ($\sigma = 1$) или определять
из первого момента уравнения энергии (добавив четвёртое ОДУ).
Начинаем с фиксированного $\sigma = 1$.


%=============================================================================
\section{Система ОДУ}\label{sec:ode}
%=============================================================================

\subsection{Свободные параметры}

Профили содержат зависящие от времени параметры:
\begin{center}
\begin{tabular}{cl}
  $r_b(t)$       & радиус зоны разгазирования, \\
  $\gamma(t)$    & показатель формы давления
                    ($\gamma = \gamma(r_b)$
                    из~\eqref{eq:gamma_constraint}), \\
  $S_{g,w}(t)$   & газонасыщенность на забое, \\
  $\Delta T_w(t)$ & отклонение температуры на забое от $T_0$.
\end{tabular}
\end{center}

Показатели формы $\beta$ и $\sigma$ на первом этапе фиксированы.
Остаются \textbf{три свободных параметра}:
$r_b(t)$, $S_{g,w}(t)$, $\Delta T_w(t)$.

\subsection{Уравнения для параметров}

Подставляя профили~\eqref{eq:p_outer}--\eqref{eq:T_profile}
в три интегральных соотношения~\eqref{eq:moment0},
\eqref{eq:Sg_moment0}, \eqref{eq:T_moment0},
получаем систему трёх ОДУ:

\paragraph{1. Из нулевого момента давления~\eqref{eq:moment0}:}
\begin{equation}\label{eq:ode1}
  \mathcal{A}_{11}\,\dot{r}_b
  + \mathcal{A}_{12}\,\dot{S}_{g,w}
  + \mathcal{A}_{13}\,\dot{\Delta T}_w
  = \mathcal{F}_1.
\end{equation}

\paragraph{2. Из нулевого момента насыщенности~\eqref{eq:Sg_moment0}:}
\begin{equation}\label{eq:ode2}
  \mathcal{A}_{21}\,\dot{r}_b
  + \mathcal{A}_{22}\,\dot{S}_{g,w}
  + \mathcal{A}_{23}\,\dot{\Delta T}_w
  = \mathcal{F}_2.
\end{equation}

\paragraph{3. Из нулевого момента энергии~\eqref{eq:T_moment0}:}
\begin{equation}\label{eq:ode3}
  \mathcal{A}_{31}\,\dot{r}_b
  + \mathcal{A}_{32}\,\dot{S}_{g,w}
  + \mathcal{A}_{33}\,\dot{\Delta T}_w
  = \mathcal{F}_3.
\end{equation}

Матричная форма:
\begin{equation}\label{eq:ode_system}
  \boxed{
    \begin{pmatrix}
      \mathcal{A}_{11} & \mathcal{A}_{12} & \mathcal{A}_{13} \\
      \mathcal{A}_{21} & \mathcal{A}_{22} & \mathcal{A}_{23} \\
      \mathcal{A}_{31} & \mathcal{A}_{32} & \mathcal{A}_{33}
    \end{pmatrix}
    \begin{pmatrix}
      \dot{r}_b \\ \dot{S}_{g,w} \\ \dot{\Delta T}_w
    \end{pmatrix}
    =
    \begin{pmatrix}
      \mathcal{F}_1 \\ \mathcal{F}_2 \\ \mathcal{F}_3
    \end{pmatrix}.
  }
\end{equation}


%=============================================================================
\section*{Литература}
%=============================================================================

\begin{thebibliography}{9}

\bibitem{barenblatt1954}
  Баренблатт~Г.\,И. О некоторых приближённых методах в теории
  одномерной неустановившейся фильтрации жидкости при упругом
  режиме~// Изв. АН СССР, ОТН. --- 1954. --- №\,9. --- С.\,35--49.

\bibitem{BER1984}
  Баренблатт~Г.\,И., Ентов~В.\,М., Рыжик~В.\,М.
  Движение жидкостей и газов в природных пластах. ---
  М.: Недра, 1984. --- 211~с.

\bibitem{mirz2004}
  Мирзаджанзаде~А.\,Х., Хасанов~М.\,М., Бахтизин~Р.\,Н.
  Моделирование процессов нефтегазодобычи. ---
  М.--Ижевск: Институт компьютерных исследований, 2004. ---
  С.\,212--218 (§\,3.7).

\bibitem{pirverdyan1961}
  Пирвердян~А.\,М. Об одном способе оценок приближённых решений
  уравнений нестационарной фильтрации нефти и газа~//
  ПММ. --- 1961. --- Т.\,25, вып.\,4.

\bibitem{entov1965}
  Ентов~В.\,М. Теоремы сравнения для уравнений нестационарной
  фильтрации~// ПММ. --- 1965. --- Т.\,29, вып.\,11.

\end{thebibliography}

\end{document}
